\chapter{การประเมินผล}

บทนี้นำเสนอแนวทางการประเมินผลของระบบที่พัฒนาขึ้น โดยมีวัตถุประสงค์เพื่อวิเคราะห์ประสิทธิภาพของเกมในด้านต่าง ๆ ทั้งด้าน gameplay ระบบปัญญาประดิษฐ์และประสบการณ์ผู้ใช้

\section{เป้าหมายการประเมิน}

การประเมินผลของโครงงานนี้มีวัตถุประสงค์เพื่อยืนยันว่า ระบบที่พัฒนาขึ้นสามารถทำงานได้ตามที่ออกแบบไว้และสร้างประสบการณ์การเล่นที่เหมาะสม โดยกำหนดเป้าหมายการประเมินออกเป็น 3 ด้านหลัก ดังนี้

\subsection{ด้านการทำงานของระบบ (System Functionality)}

เพื่อตรวจสอบว่าระบบหลักของเกมสามารถทำงานได้อย่างถูกต้องและสอดคล้องกับการออกแบบ ได้แก่

\begin{itemize}
\item ระบบการต่อสู้ (Combat System)
\item ระบบการทำอาหาร (Food System)
\item ระบบการแข่งขัน (Competition System)
\item ระบบพฤติกรรม (Behavior System)
\end{itemize}

เกณฑ์การประเมิน: ระบบสามารถทำงานได้ครบตามฟังก์ชันที่กำหนด และไม่มีข้อผิดพลาดที่ส่งผลต่อการเล่น

\subsection{ด้านความถูกต้องของแบบจำลอง (Model Validity)}

เพื่อประเมินความเหมาะสมของแบบจำลองทางคณิตศาสตร์และอัลกอริทึมที่ใช้ในระบบ โดยเฉพาะ

\begin{itemize}
\item ระบบให้คะแนนที่ใช้ Manhattan Distance
\item พฤติกรรมของตัวละคร ในแต่ละประเภท
\end{itemize}

เกณฑ์การประเมิน: ผลลัพธ์ของระบบมีความสมเหตุสมผล สอดคล้องกับการออกแบบ และสามารถ
\par\vspace{0.5em}\noindent
อธิบายได้ในเชิงตรรกะ

\subsection{ด้านประสบการณ์ผู้เล่น (Player Experience)}

เพื่อประเมินว่าผู้เล่นสามารถเข้าใจระบบ และได้รับประสบการณ์การเล่นที่ดีโดยพิจารณาจาก

\begin{itemize}
\item ความเข้าใจระบบ (Usability)
\item ความสนุกในการเล่น (Enjoyment)
\item ความต้องการเล่นซ้ำ (Replayability)
\end{itemize}

\par\vspace{0.5em}\noindent
เกณฑ์การประเมิน: ผู้เล่นสามารถเรียนรู้ระบบได้โดยไม่ซับซ้อน และมีแนวโน้มที่จะกลับมาเล่นซ้ำ

\par\vspace{0.5em}\noindent
โดยสรุป เป้าหมายของการประเมินในโครงงานนี้มุ่งเน้นการตรวจสอบทั้งในด้านการทำงานของระบบ ความถูกต้องของแบบจำลอง และประสบการณ์ของผู้เล่น เพื่อให้มั่นใจว่าเกมที่พัฒนาขึ้นมีคุณภาพและสอดคล้องกับวัตถุประสงค์ของโครงงาน

\section{วิธีการประเมิน (Evaluation Methodology)}

เพื่อให้การประเมินผลสอดคล้องกับเป้าหมายที่กำหนดไว้ในหัวข้อ 4.1 โครงงานนี้ได้ออกแบบวิธีการประเมินโดยใช้ทั้งแนวทางเชิงปริมาณ และเชิงคุณภาพ เพื่อให้ครอบคลุมทั้งด้านการทำงานของระบบและ
\par\vspace{0.5em}\noindent
ประสบการณ์ของผู้เล่น


\subsection{การทดสอบการทำงานของระบบ (Functional Testing)}

ใช้เพื่อตรวจสอบว่าระบบหลักของเกมสามารถทำงานได้อย่างถูกต้องตามที่ออกแบบ โดยครอบคลุมระบบ
\par\vspace{0.5em}\noindent
ต่าง ๆ ได้แก่

\begin{itemize}
\item ระบบการต่อสู้
\item ระบบการทำอาหาร
\item ระบบการแข่งขัน
\item ระบบปัญญาประดิษฐ์
\end{itemize}

การทดสอบจะดำเนินการโดยการทดลองใช้งานในสถานการณ์ต่าง ๆ (test scenarios) เพื่อดูว่าระบบสามารถตอบสนองได้อย่างถูกต้อง และไม่มีข้อผิดพลาดที่ส่งผลต่อ gameplay

\subsection{การทดสอบโดยผู้เล่น (User Testing)}

ใช้เพื่อประเมินประสบการณ์ของผู้เล่น โดยให้กลุ่มผู้เล่นทดลองเล่นเกม และสังเกตพฤติกรรมระหว่างการเล่น เช่น

\begin{itemize}
\item ระยะเวลาในการเรียนรู้ระบบ
\item ความสามารถในการเข้าใจ gameplay
\item การตอบสนองต่อความท้าทายของเกม
\end{itemize}

ข้อมูลที่ได้จะนำไปใช้วิเคราะห์ในด้าน usability และ player experience

\subsection{การเก็บข้อมูลแบบสอบถาม (Questionnaire)}

หลังจากการทดลองเล่น ผู้เล่นจะตอบแบบสอบถามเพื่อประเมินความรู้สึกและความคิดเห็น โดยครอบคลุมหัวข้อ เช่น

\begin{itemize}
\item ความสนุกของเกม
\item ความเข้าใจระบบ
\item ความเหมาะสมของความยาก
\item ความต้องการเล่นซ้ำ
\end{itemize}

ข้อมูลที่ได้จะถูกนำมาวิเคราะห์ในเชิงคุณภาพและเชิงปริมาณ

\subsection{การวิเคราะห์ระบบ (System Analysis)}

ใช้เพื่อตรวจสอบความถูกต้องของระบบที่ใช้แบบจำลองทางคณิตศาสตร์และ AI เช่น

\begin{itemize}
\item ตรวจสอบผลลัพธ์ของระบบให้คะแนน (Scoring System)
\item วิเคราะห์พฤติกรรมของตัวละคร ว่าตรงตามที่ออกแบบไว้
\end{itemize}

\section{ตัวชี้วัดการประเมินผล (Evaluation Metrics)}

เพื่อให้การประเมินผลมีความชัดเจนและสามารถวิเคราะห์ได้ในเชิงปริมาณ โครงงานนี้ได้กำหนดตัวชี้วัด สำหรับการประเมินในแต่ละด้าน โดยแบ่งออกเป็น 3 กลุ่มหลัก ดังนี้

\subsection{ตัวชี้วัดด้าน Gameplay}

ใช้สำหรับประเมินประสิทธิภาพของการออกแบบการเล่น และความสมดุลของระบบ

\begin{itemize}
\item \textbf{Average Play Time per Session} \\
ระยะเวลาเฉลี่ยที่ผู้เล่นใช้ในแต่ละรอบการเล่น

\item \textbf{Win/Loss Ratio} \\
อัตราการชนะและแพ้ในการแข่งขัน

\item \textbf{Progression Rate} \\
ความเร็วในการปลดล็อกเนื้อหาใหม่ของผู้เล่น
\end{itemize}

\subsection{ตัวชี้วัดด้านระบบ (System Metrics)}

ใช้สำหรับประเมินความถูกต้องและประสิทธิภาพของระบบที่พัฒนา

\begin{itemize}
\item \textbf{Scoring Accuracy} \\
ความสอดคล้องของคะแนนที่ได้กับความใกล้เคียงของรสชาติ

\item \textbf{AI Response Consistency} \\
ความสม่ำเสมอของพฤติกรรม AI ในสถานการณ์เดียวกัน

\item \textbf{System Stability} \\
จำนวนข้อผิดพลาด (errors) หรือบั๊กที่เกิดขึ้นระหว่างการเล่น
\end{itemize}

\subsection{ตัวชี้วัดด้านประสบการณ์ผู้ใช้ (User Experience Metrics)}

ใช้สำหรับประเมินความพึงพอใจและความเข้าใจของผู้เล่น โดยเก็บข้อมูลผ่านแบบสอบถาม

\begin{itemize}
\item \textbf{Fun Score} \\
ระดับความสนุกในการเล่น

\item \textbf{Usability Score} \\
ความง่ายในการเรียนรู้และใช้งานระบบ

\item \textbf{Replayability Score} \\
ความต้องการเล่นซ้ำของผู้เล่น
\end{itemize}

\par\vspace{0.5em}\noindent
โดยตัวชี้วัดในส่วนนี้สามารถใช้มาตราส่วนคะแนน (Likert Scale) เช่น 1--5 เพื่อวิเคราะห์ผลลัพธ์

\par\vspace{0.5em}\noindent
โดยสรุป การกำหนดตัวชี้วัดดังกล่าวช่วยให้สามารถวิเคราะห์ผลการประเมินได้อย่างเป็นระบบ และสามารถเปรียบเทียบผลลัพธ์ในแต่ละด้านได้อย่างชัดเจน


\section{การออกแบบการทดลอง (Experimental Setup)}

เพื่อให้การประเมินผลสามารถดำเนินการได้อย่างเป็นระบบ โครงงานนี้ได้ออกแบบขั้นตอนการทดลอง (experimental setup) สำหรับใช้ในการทดสอบเกม โดยครอบคลุมทั้งด้านการทำงานของระบบและ
\par\vspace{0.5em}\noindent
ประสบการณ์ของผู้เล่น

\subsection{กลุ่มผู้ทดลอง (Participants)}

การทดลองจะใช้กลุ่มผู้เล่นตัวอย่างจำนวนหนึ่ง โดยแบ่งออกเป็นกลุ่มตามระดับประสบการณ์ในการเล่นเกม ได้แก่:

\begin{itemize}
\item ผู้เล่นทั่วไป (Casual Players)
\item ผู้เล่นที่มีประสบการณ์ (Experienced Players)
\end{itemize}

การแบ่งกลุ่มนี้ช่วยให้สามารถวิเคราะห์ผลกระทบของระดับประสบการณ์ต่อการเล่นและการเรียนรู้ระบบ

\subsection{ขั้นตอนการทดลอง (Procedure)}

การทดลองถูกออกแบบเป็นลำดับขั้นตอนดังนี้:

\begin{enumerate}
\item ผู้เล่นได้รับคำแนะนำเกี่ยวกับระบบพื้นฐานของเกม
\item ผู้เล่นทดลองเล่นเกมในระยะเวลาที่กำหนด
\item ผู้เล่นทำกิจกรรมหลัก เช่น การต่อสู้ การทำอาหาร และการแข่งขัน
\item เก็บข้อมูลระหว่างการเล่น เช่น เวลา คะแนน และพฤติกรรม
\item ผู้เล่นตอบแบบสอบถามหลังการทดลอง
\end{enumerate}

\subsection{ตัวแปรที่ควบคุม (Controlled Variables)}

เพื่อให้ผลการทดลองมีความน่าเชื่อถือ ได้มีการกำหนดตัวแปรที่ควบคุม ได้แก่:

\begin{itemize}
\item ระยะเวลาในการทดลอง
\item ระดับความยากของเกม
\item เงื่อนไขของการแข่งขัน
\item สภาพแวดล้อมในการเล่น
\end{itemize}

\subsection{ข้อมูลที่เก็บรวบรวม (Data Collection)}

ข้อมูลที่เก็บในการทดลองแบ่งออกเป็น 2 ประเภท ได้แก่

\begin{itemize}
\item \textbf{ข้อมูลเชิงปริมาณ (Quantitative Data)}
\begin{itemize}
\item คะแนนจากการแข่งขัน
\item ระยะเวลาในการเล่น
\item อัตราการชนะ/แพ้
\end{itemize}

\item \textbf{ข้อมูลเชิงคุณภาพ (Qualitative Data)}
\begin{itemize}
\item ความคิดเห็นจากแบบสอบถาม
\item ความพึงพอใจของผู้เล่น
\item ข้อเสนอแนะในการปรับปรุง
\end{itemize}
\end{itemize}

โดยสรุป การออกแบบการทดลองนี้ช่วยให้สามารถประเมินผลระบบได้อย่างเป็นระบบ และสามารถนำไปใช้ในการทดสอบจริงในอนาคตได้อย่างมีประสิทธิภาพ


\section{ข้อจำกัดของการประเมิน (Limitations)}

แม้ว่าโครงงานนี้ได้มีการออกแบบแนวทางการประเมินผลอย่างเป็นระบบ แต่ยังมีข้อจำกัดบางประการที่อาจส่งผลต่อความสมบูรณ์ของการประเมิน ดังนี้

\subsection{การยังไม่ได้ดำเนินการทดลองจริง}

เนื่องจากโครงงานอยู่ในขั้นตอนการพัฒนา การประเมินผลเชิงทดลองยังไม่ได้ดำเนินการจริง ทำให้ยังไม่มี
\par\vspace{0.5em}\noindent
ข้อมูลเชิงประจักษ์ (empirical data) เพื่อยืนยันผลลัพธ์ของระบบ

\subsection{ขนาดของกลุ่มตัวอย่าง}

ในการทดลองในอนาคต อาจมีข้อจำกัดด้านจำนวนผู้เข้าร่วม ซึ่งอาจส่งผลต่อความน่าเชื่อถือของผลการ 
\par\vspace{0.5em}\noindent
วิเคราะห์โดยเฉพาะในด้านประสบการณ์ผู้ใช้

\subsection{ความหลากหลายของผู้เล่น}

ผู้เล่นที่เข้าร่วมการทดลองอาจมีพื้นฐานหรือประสบการณ์ที่แตกต่างกัน ซึ่งอาจส่งผลต่อผลลัพธ์ เช่น ความเข้าใจระบบ หรือระดับความสนุกในการเล่น

\subsection{ข้อจำกัดด้านเวลาและทรัพยากร}

ระยะเวลาในการพัฒนาและทดสอบอาจมีจำกัด ทำให้ไม่สามารถทดสอบระบบในทุกสถานการณ์
\par\vspace{0.5em}\noindent
หรือปรับปรุงระบบได้อย่างเต็มที่

\subsection{แนวทางการพัฒนาในอนาคต}

ในอนาคต ควรมีการดำเนินการทดลองจริงกับกลุ่มผู้เล่น และเก็บข้อมูลเชิงปริมาณอย่างเป็นระบบ เพื่อ
\par\vspace{0.5em}\noindent
ปรับปรุงและเพิ่มความแม่นยำของระบบ โดยเฉพาะในส่วนของ Behavior และระบบการให้คะแนน

\par\vspace{1em}\noindent
โดยสรุป แม้จะมีข้อจำกัดบางประการ แต่โครงงานนี้ได้มีการวางแผนการประเมินผลไว้อย่างชัดเจน 
\par\vspace{0.5em}\noindent
ซึ่งสามารถนำไปพัฒนาและปรับปรุงต่อได้ในอนาคต